% Generated draft: requires mathematical and source review.
\documentclass[11pt]{article}
\usepackage[margin=25mm]{geometry}
\usepackage{fontspec}
\setmainfont{Latin Modern Roman}
\usepackage{amsmath,amssymb,mathtools,unicode-math}
\setmathfont{Latin Modern Math}
\usepackage{enumitem,needspace,xurl,hyperref}
\setlength{\parindent}{0pt}
\setlength{\parskip}{5pt}
\setlength{\emergencystretch}{3em}
\newcommand{\N}{\mathbb{N}}
\newcommand{\Z}{\mathbb{Z}}
\newcommand{\Q}{\mathbb{Q}}
\newcommand{\R}{\mathbb{R}}
\newcommand{\C}{\mathbb{C}}
\newcommand{\F}{\mathbb{F}}
\newcommand{\A}{\mathbb{A}}
\newcommand{\PP}{\mathbb{P}}
\newcommand{\E}{\mathbb{E}}
\newcommand{\eps}{\varepsilon}
\newcommand{\dd}{\,\mathrm{d}}
\begin{document}
\section*{UnsolvedMath ID: 700559}
\textbf{Source problem code:} catalogue entry
\subsection{Definitions and mathematical statement}
Let \(G\) be a finite group. Its \(3\)-rank is
\[
r_3(G)=\max\{d\in\Z_{\geq 0}:G\text{ contains a subgroup isomorphic to }(\Z/3\Z)^d\}.
\]
For a nontrivial \(2\)-subgroup \(Q\leq G\), write
\[
N_G(Q)=\{g\in G:gQg^{-1}=Q\}.
\]
The supplied record does not define ``\(2\)-local \(3\)-rank.'' Under the convention that \(2\)-local subgroups are normalizers of nontrivial \(2\)-subgroups, define
\[
r_{2,3}(G)=\max\bigl(\{0\}\cup\{r_3(N_G(Q)):1\neq Q\leq G,\ |Q|=2^a\text{ for some integer }a\geq1\}\bigr).
\]
The added value \(0\) specifies the convention when there are no nontrivial \(2\)-subgroups. This interpretation must be checked against the original source.

\textbf{Part (b), the remaining classification request.} Subject to that interpretation, describe, up to isomorphism, all finite groups \(G\) such that
\[
r_{2,3}(G)=1
\quad\text{and}\quad
P\text{ is noncyclic for }P\in\operatorname{Syl}_3(G).
\]
Here \(\operatorname{Syl}_3(G)\) denotes the set of Sylow \(3\)-subgroups of \(G\). Since these subgroups are conjugate, the noncyclicity condition holds for one if and only if it holds for every Sylow \(3\)-subgroup.

\textbf{Part (a), included as supplied archived material.} Describe, up to isomorphism, all finite groups \(G\) satisfying
\[
r_{2,3}(G)=1\quad\text{and}\quad r_3(G)\geq3.
\]
The supplied record labels part (a) as archived solved and part (b) as remaining in the main section. These are supplied status qualifications, not independently verified mathematical conclusions.
\subsection{Short English statement}
Part (b) asks for a classification of finite groups of \(2\)-local \(3\)-rank exactly \(1\) with noncyclic Sylow \(3\)-subgroups. The definition of local rank used here requires source confirmation. Part (a), concerning groups of \(3\)-rank at least \(3\), is retained as supplied archived solved material.
\subsection{Sources}
\begin{itemize}
\item Local UnsolvedMath registry, permanent ID \(700559\), source code \texttt{catalogue entry}, title \emph{Kourovka Notebook Problem 9.26}; supplied upstream revision \texttt{d371ee917ac9647fbb06cd6fc529deec6f3dda82}.
\item \emph{The Kourovka Notebook (Unsolved Problems in Group Theory)}, problem 9.26, proposed by A. A. Makhn\"ev in the 9th issue (1984); the supplied record reports an update in the 21st issue (2026), arXiv:1401.0300v48. The registry describes the collection as Issues 1--21 (1965--2026), edited by E. I. Khukhro and V. D. Mazurov.
\item Supplied Archive of Solved Problems attribution: A. A. Makhn\"ev, \emph{Siber. Math. J.} \textbf{29} (1988), 951--959. The archive entry says the groups are described there; the supplied status review associates the archived solved material with part (a). No article title is supplied.
\item The supplied status review, checked on \texttt{2026-10-16}, labels the record \texttt{partially\_solved}, with solved part(s) in the archive and remaining part(s) in the main section; its evidence field is \texttt{published}. This is supplied dated curation, not independently verified source or literature assessment.
\end{itemize}
\end{document}
